%----------------------------------------------------------------
% Poglavje (Chapter) 1
%----------------------------------------------------------------
\chapter{Uvod}
\label{ch:uvod}

Datoteka {\tt magistrska\_naloga.tex} na kratko opisuje, kako se pisanja magistrskega dela lotimo z uporabo
programskega pateka \LaTeX. V tem dokumentu bomo predstavili nekaj njegovih prednosti in hib. Kar se slednjih
tiče, mi pride na misel ena sama. Ko se srečamo z njim, nam izgleda kot kislo jabolko, nismo prepričani, da
bi želeli vanj ugrizniti. Lahko pa z njim pripravimo odličen zavitek ali pa pridemo na okus.

V Poglavju~\ref{ch:uvod} bomo na hitro spoznali besedilne konstrukte kot so izreki, enačbe in dokazi. Naučili
se bomo, kako se na njih sklicujemo. V Poglavju~\ref{ch:sklicevanje} se bomo srečali s sklicevanjem na
besedilne konstrukte. Poglavje~\ref{ch:plovke} bo predstavilo vključevanje plovk: slik in tabel. V
Poglavju~\ref{ch3} se bomo srečali s sklicevanjem na literaturo.
Sledil bo samo še zaključek.

Bodite pozorni, da se v glavni mapi nahajata še datoteki \verb+ declaration.tex+ in \verb+ izjava.tex+. Ti
datoteki se ločeno prevedeta, ju podpišete in oddate v referat ločeno od magistrske naloge.